\begin{examplebox}
	\textbf{\textcolor{CExample}{例 1（基础）}}

	\textbf{\textcolor{CExample}{题目：}}
	如图，$l_1\parallel l_2\parallel l_3$，直线 $a$ 分别交 $l_1,l_2,l_3$ 于点 $A,B,C$，直线 $b$ 分别交 $l_1,l_2,l_3$ 于点 $D,E,F$。已知 $AB=2$，$BC=3$，$DE=4$，求 $EF$ 和 $DF$ 的长。

	\textbf{\textcolor{CExample}{解题步骤：}}
	\begin{description}[style=nextline, leftmargin=2.8em, labelsep=0.5em, itemsep=0.45em]
		\item[\textbf{第一步（明确对应关系）：}] 根据平行线顺序，$AB$ 对应 $DE$，$BC$ 对应 $EF$，$AC$ 对应 $DF$。
			\par\textbf{理由：}
			定理指出对应线段成比例。
		\item[\textbf{第二步（求 EF）：}] 由 $\dfrac{AB}{BC}=\dfrac{DE}{EF}$ 代入 $AB=2,BC=3,DE=4$，得 $\dfrac{2}{3}=\dfrac{4}{EF}$，解得 $EF=6$。
			\par\textbf{理由：}
			交叉相乘即可。
		\item[\textbf{第三步（求 DF）：}] 由 $\dfrac{AB}{AC}=\dfrac{DE}{DF}$，先算 $AC=AB+BC=5$，代入得 $\dfrac{2}{5}=\dfrac{4}{DF}$，解得 $DF=10$。
			\par\textbf{理由：}
			交叉相乘。
	\end{description}

	\textbf{\textcolor{CExample}{关键结论：}}
	\textbf{$\boxed{EF=6,\ DF=10}$}

	\medskip
	\textbf{\textcolor{CExample}{例 2（稍变形）}}

	\textbf{\textcolor{CExample}{题目：}}
	在 $\triangle ABC$ 中，$DE\parallel BC$，交 $AB$ 于 $D$，$AC$ 于 $E$。已知 $AD=2$，$DB=3$，$AC=10$，求 $AE$ 和 $EC$ 的长。

	\textbf{\textcolor{CExample}{解题步骤：}}
	\begin{description}[style=nextline, leftmargin=2.8em, labelsep=0.5em, itemsep=0.45em]
		\item[\textbf{第一步（识别平行线组）：}] 图中 $DE\parallel BC$，过 $A$ 作 $BC$ 的平行线（即与 $DE$ 平行），从而有三条平行线，根据平行线分线段成比例的常见变形，得 $\dfrac{AD}{DB}=\dfrac{AE}{EC}$。
			\par\textbf{理由：}
			平行线组截 $AB$ 和 $AC$。
		\item[\textbf{第二步（设未知数）：}] 设 $AE=2k$，则 $EC=3k$，由 $AE+EC=AC=10$ 得 $2k+3k=10$，解得 $k=2$。
			\par\textbf{理由：}
			利用比例式设参数。
		\item[\textbf{第三步（求结果）：}] $AE=2k=4$，$EC=3k=6$。
			\par\textbf{理由：}
			回代。
	\end{description}

	\textbf{\textcolor{CExample}{关键结论：}}
	\textbf{$\boxed{AE=4,\ EC=6}$}

	\medskip
	\textbf{\textcolor{CExample}{例 3（综合应用）}}

	\textbf{\textcolor{CExample}{题目：}}
	如图，$l_1\parallel l_2\parallel l_3$，直线 $a$ 分别交 $l_1,l_2,l_3$ 于 $A,B,C$，直线 $b$ 分别交 $l_1,l_2,l_3$ 于 $D,E,F$。已知 $AB=3$，$BC=5$，$DE+EF=16$，求 $DE$ 和 $EF$。

	\textbf{\textcolor{CExample}{解题步骤：}}
	\begin{description}[style=nextline, leftmargin=2.8em, labelsep=0.5em, itemsep=0.45em]
		\item[\textbf{第一步（列比例式）：}] 由 $\dfrac{AB}{BC}=\dfrac{DE}{EF}$ 代入 $AB=3,BC=5$ 得 $\dfrac{3}{5}=\dfrac{DE}{EF}$，即 $DE=\dfrac{3}{5}EF$。
			\par\textbf{理由：}
			平行线分线段成比例定理的第一个比例式。
		\item[\textbf{第二步（利用和条件）：}] 由 $DE+EF=16$，代入 $DE=\dfrac{3}{5}EF$，得 $\dfrac{3}{5}EF+EF=16$，即 $\dfrac{8}{5}EF=16$，解得 $EF=10$。
			\par\textbf{理由：}
			代数运算。
		\item[\textbf{第三步（求 DE）：}] $DE=\dfrac{3}{5}\times10=6$。
			\par\textbf{理由：}
			回代。
	\end{description}

	\textbf{\textcolor{CExample}{关键结论：}}
	\textbf{$\boxed{DE=6,\ EF=10}$}
\end{examplebox}
